% path to figures directory
\graphicspath{{img/chapter_1/}}

\chapter{Introduction}
\label{chapter:introduction}

\begin{flushleft}
  \textit{The following is a general introduction to the background topics
    referred to and assumed in subsequent chapters. This includes a review of
    popular theories of neutrino mass, the current status of
    neutrino-oscillation parameters, a general introduction to Effective Field
    Theory (EFT), the experimental situation relevant to the flavour anomalies,
    and topics peripheral to all of these.}
\end{flushleft}

\section{The Standard Model and neutrinos}

Laboratory experiments to date have firmly established the predictive power of
the Standard Model (SM) of particle physics, a combined theory of the
electroweak and strong interactions described by the gauge group
$G_{\text{SM}} = \mathrm{SU}(3)_{c} \otimes \mathrm{SU}(2)_{L} \otimes \mathrm{U}(1)_{Y}$.
It is a model whose probes and predictions span at least 33 orders of
magnitude\footnote{The interval given is from the distance scales probed at the
  LHC (roughly $\SI{e-17}{\cm}$) to the size of the solar system (roughly
  \SI{e16}{\cm}).} with varying degrees of precision, and these are consistent
with almost all known experiments. Although it displays a number of arbitrary
features, the dynamics of the theory are mostly fixed by the fundamental
principles of gauge theory and Lorentz invariance. Most of this arbitrariness
resides in the matter sector of the theory, whose properties (masses, coupling
constants, quantum numbers, \textit{etc.}) are not predicted, but are instead
motivated on phenomenological grounds. We show the fields of the SM and their
defining properties in Table~\ref{tab:sm-fields}, according to the mathematical
conventions of Appendix~\ref{chapter:notation}.
\begin{table}[t]
  \centering
  \bgroup
  \def\arraystretch{1.3}% 1 is the default
  \begin{tabular}{ccc}
    \toprule
    Field                        & $\mathrm{SU}(3)_{c} \otimes \mathrm{SU}(2)_{L} \otimes \mathrm{U}(1)_{Y}$ & $\mathrm{SU}(2)_{+} \otimes \mathrm{SU}(2)_{-}$ \\
    \midrule
    $Q^{\alpha a i}$             & $(\mathbf{3}, \mathbf{2}, \tfrac{1}{6})$                                  & $(\mathbf{2}, \mathbf{1})$                      \\
    $L^{\alpha i}$               & $(\mathbf{1}, \mathbf{2}, -\tfrac{1}{2})$                                 & $(\mathbf{2}, \mathbf{1})$                      \\
    $\bar{u}^{\alpha}_a$                  & $(\bar{\mathbf{3}}, \mathbf{1}, -\tfrac{2}{3})$                           & $(\mathbf{2}, \mathbf{1})$                      \\
    $\bar{d}^{\alpha}_a$                  & $(\bar{\mathbf{3}}, \mathbf{1}, \tfrac{1}{3})$                            & $(\mathbf{2}, \mathbf{1})$                      \\
    $\bar{e}^{\alpha}$                    & $(\mathbf{1}, \mathbf{1}, 1)$                                             & $(\mathbf{2}, \mathbf{1})$                      \\
    $(G_{\alpha \beta})^a_{\ b}$ & $(\mathbf{8}, \mathbf{1}, 0)$                                             & $(\mathbf{3}, \mathbf{1})$                      \\
    $(W_{\alpha \beta})^i_{\ j}$ & $(\mathbf{1}, \mathbf{3}, 0)$                                             & $(\mathbf{3}, \mathbf{1})$                      \\
    $B_{\alpha \beta}$           & $(\mathbf{1}, \mathbf{1}, 0)$                                             & $(\mathbf{3}, \mathbf{1})$                      \\
    $H^{i}$                      & $(\mathbf{1}, \mathbf{2}, \tfrac{1}{2})$                                  & $(\mathbf{1}, \mathbf{1})$                      \\
    \bottomrule
  \end{tabular}
  \egroup
  \caption{The SM fields and their transformation properties under the SM gauge
    group $G_{\text{SM}}$ and the Lorentz group written as
    $\mathrm{SU}(2)_{+} \otimes \mathrm{SU}(2)_{-}$. The final unbolded number
    in the 3-tuples of the $G_{\text{SM}}$ column represents the
    $\mathrm{U}(1)_Y$ charge of the field, normalised such that $Q = I_{3} + Y$.
    For the fermions a generational index has been suppressed. See
    Appendix~\ref{chapter:notation} for details about the mathematical notation
    used here and throughout this work.}
  \label{tab:sm-fields}
\end{table}

The SM inherits the experimental success of the
$\mathrm{SU}(2) \otimes \mathrm{U}(1)$ theory of the weak interactions, first
proposed by Glashow~\cite{Glashow:1961tr} in 1961 as a possible underlying
structure for Fermi's theory of beta decay. Before the end of the same decade,
Weinberg~\cite{Weinberg:1967tq} and Salam~\cite{Salam:1968rm} had constructed
the modern theory of leptons based on the spontaneous breaking of
$\mathrm{SU}(2)_{L} \otimes \mathrm{U}(1)_{Y}$ to the electromagnetic symmetry.
Interestingly, it seems that these seminal papers went mostly unnoticed (see
Fig.~\ref{fig:weinberg-citations}) until the early 1970s, when 't~Hooft proved
the renormalisability of spontaneously broken gauge
theories~\cite{tHooft:1971akt} as a graduate student working under the
supervision of Veltman. By the mid 1970s the framework had been extended to
include the quarks~\cite{Glashow:1970gm} and the unbroken chromodynamic group,
which was successfully shown to reproduce the Bjorken scaling seen in
deep-inelastic-scattering experiments through asymptotic
freedom~\cite{Gross:1973id}.
\begin{figure}[t]
  \centering
  \includegraphics[width=0.9\linewidth]{weinberg_citations}
  \caption{The cumulative citation graph for a selection of papers presenting
    foundational results relevant to the SM. Weinberg's seminal paper `A model
    of leptons' (1967) saw an explosion of citations following 't~Hooft's work
    on the renormalisability of gauge theories (1971).}
  \label{fig:weinberg-citations}
\end{figure}

Despite its successes, the SM cannot be the complete theory of fundamental
particles and their interactions. It does not explain phenomena such as the
baryon asymmetry present, and the particle spectrum contains no viable candidate
for dark matter. The SM cannot explain why the electric dipole moment of the
neutron is so small, why there are three generations of matter or, notably in
our case, the origin of neutrino oscillations and the implied small but non-zero
neutrino masses.

\section{Massive neutrinos in experiment and theory}

The minimal SM predicts massless neutrinos, a prediction that today sits in
contradiction to a wealth of empirical evidence. This evidence could in
principle have come from many kinds of experiments, but currently only neutrino
oscillations provide strong signs that the masses are non-zero. Below we discuss
the phenomenon of neutrino oscillations in the context of the outstandingly
successful three-flavour mixing paradigm. We then move on to other probes of
neutrino masses, which currently only provide limits on the mass scale. On the
theory side, we summarise some popular and motivated extensions of the SM that
accommodate massive neutrinos, placing particular emphasis on the direction we
have followed in the novel work presented in this thesis. This includes an
overview of tree- and loop-level models of Majorana neutrino mass.

\subsection{Neutrino oscillations}

The neutrino flavour eigenstates $\breve{\nu}_{i}$ =
$(\nu_{e}, \nu_{\mu}, \nu_{\tau})$ are defined as the states that couple at
charged-current interaction vertices with the corresponding charged lepton.
These are the states in which the neutrinos are almost always produced in
experiments, and certainly always measured. If neutrinos are massive there is no
reason to expect these to coincide with the mass eigenstates
$\nu_{i} = (\nu_{1}, \nu_{2}, \nu_{3})$. In general, the flavour eigenstates
will be an admixture of the propagating fields
\begin{equation}
  \label{eq:neutrino-mixing}
  \breve{\nu}_{i} = U_{i}^{\ j} \nu_{j} \ ,
\end{equation}
where the $U_{i}^{\ j}$ are elements of the unitary
Pontecorvo--Maki--Nakagawa--Sakata (PMNS) neutrino mixing
matrix~\cite{Pontecorvo:1957cp, Maki:1962mu}. The PMNS matrix is defined such
that it diagonalises the neutrino mass matrix:
\begin{equation}
  \mathbf{U}^{\dagger} \mathbf{m}_{\nu} \mathbf{U}^{*} = \mathrm{diag}(m_{1}, m_{2}, m_{3}),
\end{equation}
where the $m_i$ are the neutrino masses. Being a $3 \times 3$ unitary matrix,
$\mathbf{U}$ is in general parametrised by three mixing angles and six phases.
Not all of the phases are physical, since the neutrino and charged-lepton fields
can be redefined in such a way that five of the phases are removed in the case
of Dirac neutrinos. In the presence of a Majorana mass term, only the charged
leptons can be rephased. This leaves three physical phases with the two
additional ones termed \textit{Majorana phases}. In general
\begin{equation}
  \label{eq:pmns}
  \mathbf{U} =
  \begin{bmatrix}
    c_{12}c_{13} & s_{12} c_{13} & s_{13}e^{-i\delta_\text{CP}} \\
    -s_{12}c_{23} - c_{12}s_{23}s_{13}e^{i\delta_\text{CP}} & c_{12}c_{23} - s_{12}s_{23}s_{13}e^{i\delta_\text{CP}} & s_{23}c_{13}\\
    s_{12}s_{23} - c_{12}c_{23}s_{13}e^{i\delta_\text{CP}} & -c_{12}s_{23} - s_{12}c_{23}s_{13}e^{i\delta_\text{CP}} & c_{23}c_{13}
  \end{bmatrix}
  \mathbf{P} \ ,
\end{equation}
where $c_{ij} = \cos \theta_{ij}$, $s_{ij} = \sin \theta_{ij}$ and
\begin{equation}
  \mathbf{P} =
  \begin{cases}
    \mathrm{diag}(e^{i \alpha_{1}}, e^{i \alpha_{2}}, 1) & \text{for Majorana neutrinos} \\
    \mathbf{1}_{3 \times 3} & \text{for Dirac neutrinos} \\
  \end{cases} \ .
\end{equation}
The phase $\delta_{\text{CP}}$ is often called the \textit{Dirac phase}, while
$\alpha_{1,2}$ are the Majorana phases discussed above.

Neutrino oscillation experiments typically involve the production of neutrino
flavour states from charged-current processes, \textit{e.g.} leptonic pion
decays. Each mass eigenstate evolves in time independently according to the
Schr\"{o}dinger equation: \textit{i.e.}
$| \nu_{i} (t) \rangle = \exp{(-i E_{i} t)} | \nu_{i} (0) \rangle$, for evolution
\textit{in vacuo}. This alters the initial superposition away from being a pure
flavour eigenstate:
\begin{align}
  | \breve{\nu}_{i} (t) \rangle &= \sum_{j} U_{i}^{* j} e^{-i E_{j} t} | \nu_{j} \rangle \\
                                &= \sum_{j,k} U_{i}^{* j} e^{-i E_{j} t} U_{j}^{\ k} | \breve{\nu}_{k} \rangle \ .
\end{align}
The probability of measuring a specific flavour through the charged-current
interaction then oscillates with time:
\begin{equation}
  \text{P}(\breve{\nu}_{m} \to \breve{\nu}_{n}) = | \langle \breve{\nu}_{n} | \breve{\nu}_{m} (t) \rangle |^{2} = \left| \sum_{i} U_{m}^{* i} U_{i}^{\ n} e^{-i E_{i} t} \right|^{2} \ .
\end{equation}
The expression can be expanded and the kinematic factors simplified from the
fact that the neutrinos are ultra-relativistic. We follow the usual convention
and take
$E_{i} = \sqrt{\mathbf{p}^{2} + m_{i}^{2}} \approx |\mathbf{p}| + m_{i}^{2} / (2 E)$
with $E = |\mathbf{p}|$. This gives
\begin{equation}
  \label{eq:neutrino-osc}
  \begin{aligned}
    \text{P}(\breve{\nu}_{m} \to \breve{\nu}_{n}) &= \delta_{mn} - 4 \sum_{i < j} \mathrm{Re}\left( U_{m}^{\ i} U_{m}^{* j} U_{n}^{* i} U_{n}^{\ j}  \right) \sin^{2} \frac{\Delta m_{ij}^{2} L}{4E}\\
    &\quad + 2 \sum_{i<j} \mathrm{Im}\left( U_{m}^{\ i} U_{m}^{* j} U_{n}^{* i} U_{n}^{\ j}  \right) \sin \frac{\Delta m^{2}_{ji} L}{2E} \ ,
  \end{aligned}
\end{equation}
where $\Delta m_{ij}^{2} \equiv m_{j}^{2} - m_{i}^{2}$ are the squared neutrino
mass differences and $L = ct$, sometimes called the baseline, is the approximate
distance travelled by the particles. To interpret the results of many
experiments, it is often sufficient to consider an effective two-flavour
oscillation paradigm. In this case, the neutrino-oscillation probabilities are
governed by by a single squared mass difference $\Delta m^{2}$ and a single
angle $\theta$. Interestingly, the CP-violating phase is completely absent from
the two flavour formula:
\begin{equation}
  \label{eq:two-flavour-osc}
  \text{P}(\breve{\nu}_{m} \to \breve{\nu}_{n})_{n_{f}=2} = \sin^{2} (2\theta) \sin^{2} \frac{\Delta m^{2} L}{4E} \ .
\end{equation}

From the expressions in Eqs.~\eqref{eq:neutrino-osc} and
\eqref{eq:two-flavour-osc} a number of properties of the vacuum neutrino
oscillations become clear.
\begin{enumerate}
  \item The neutrino oscillation probabilities depend on the neutrino mass
    differences, and not on the absolute mass scale. For three flavours, there
    are only two independent squared mass differences. Typically chosen to be
    $\Delta m_{21}^{2}$ and $\Delta m_{31}^{2}$, although often they are
    refereed to with the historical names $\Delta m_{\text{sol}}^{2}$ and
    $\Delta m_{\text{atm}}^{2}$, discussed in detail below.
  \item From Eq.~\eqref{eq:neutrino-osc} it is clear that the oscillations only
    occur if the neutrinos are non-degenerate and the neutrino mixing is
    non-trivial, \textit{i.e.} if $\Delta m_{ij} \neq 0$ and
    $\mathbf{U} \neq \mathbf{1}$.
  \item The PMNS matrix elements only appear in the combination
    $U_{m}^{\ i} U_{m}^{* j}$, to which the Majorana phases contained within the
    matrix $\mathbf{P}$ do not contribute. This implies that oscillation
    experiments cannot comment on the Dirac or Majorana nature\footnote{Of
    course, this is already clear from the fact that neutrino oscillations
    conserve total lepton number, despite breaking the individual familial
    lepton-number symmetries $L_{e, \mu, \tau}$.} of the neutrinos. Oscillations
    can however probe $\delta_{\text{CP}}$.
  \item In the effective two-flavour mixing paradigm, both $\theta$ and
    $\Delta m^{2}$ appear in such a way that neither the sign of $\Delta m^{2}$
    nor the octant of $\theta$ can be uniquely determined.
  \end{enumerate}
  Thus, neutrino oscillations imply that the neutrino masses of at least two of
  the mass eigenstates are non-degenerate, and therefore only one neutrino could
  potentially be massless. The largest squared mass difference can therefore be
  translated into an lower bound on the mass of the heaviest neutrino, which we
  present later with modern data.

  Historically, the effective two-flavour mixing paradigm has provided a good
  framework for interpreting early indications of neutrino oscillations.
  Specifically, the solar and atmospheric neutrino puzzles have approximate
  descriptions in terms of a two-flavour picture. Oscillations
  $\nu_{e} \to \nu_{\text{active}}$, where $\nu_{\text{active}}$ is a coherent
  superposition of $\nu_{\mu}$ and $\nu_{\tau}$, in both matter and vacuum
  account for the deficit of electron neutrinos measured from the sun, and
  $\nu_{\mu} \to \nu_{\tau}$ oscillations \textit{in vacuo} explain the shortage
  of muon neutrinos from cosmic-ray-induced production in the upper atmosphere.

  The measurement and resolution of these puzzles is an interesting and exciting
  chapter in the recent history of physics. Experiments as early as the 1960s
  had noticed a shortage of electron neutrinos coming from the sun relative to
  the predictions of solar models~\cite{RevModPhys.60.297, 1988ApJ...335..415T,
    RevModPhys.64.885, RevModPhys.67.781}, which themselves were subject to much
  uncertainty~\cite{Morrison:1992bz}. For detection there were three main
  approaches: Raymond Davis and collaborators~\cite{PhysRevLett.12.303}
  pioneered experiments that measured the solar electron-neutrino flux using
  Chlorine, the Kamiokande and later Super-Kamiokande
  collaborations~\cite{Hirata:1989zj, Hirata:1990xa} used water Cherenkov
  detectors, and the experiments GALLEX~\cite{Hampel:1998xg} and
  SAGE~\cite{Abdurashitov:1999zd} had Gallium as the detecting material. All of
  these experiments showed a deficit of solar electron neutrinos, although they
  were sensitive to neutrinos of different energies. The Sudbury Neutrino
  Observatory gave the final word on the oscillation solution to the solar
  neutrino puzzle with accurate confirmation of the electron-neutrino deficit,
  along with a measurement of the \textit{total} neutrino flux which was found
  to be in agreement with the solar models~\cite{Ahmad:2001an, Ahmad:2002jz}.

  Atmospheric neutrinos were known to come about from helicity-suppressed kaon
  and pion decays to muons and muon neutrinos. A zenith-angle and
  energy-dependent suppression in the flux of atmospheric muon neutrinos was
  measured by the Kamiokande and Irvine--Michigan--Brookhaven
  experiments~\cite{Hirata:1992ku, BeckerSzendy:1995vr} in the early 1990s, and
  after the upgrade to Super-Kamiokande the deficit was confirmed to high
  precision with results presented at the `Neutrino 1998'
  conference~\cite{Fukuda:1998mi, vonFeilitzsch:2003mh, shiozawa:2002,
    Smy:2002rz}.

  The pairs of mixing parameters associated with these two classes of
  measurements are usually dubbed
  $\theta_{\text{sol}}, \Delta m^{2}_{\text{sol}}$ and
  $\theta_{\text{atm}}, \Delta m^{2}_{\text{atm}}$. Experimental results find
  $\Delta m^{2}_{\text{sol}} \ll \Delta m^{2}_{\text{atm}}$ and that both
  $\theta_{\text{sol}}$ and $\theta_{\text{atm}}$ are large compared to any
  angles found in the CKM matrix, the quark mixing matrix. Interpreted in terms
  of three-flavour mixing, common convention identifies
  $\Delta m_{\text{sol}}^{2}$ with the squared mass difference between $\nu_{2}$
  and $\nu_{1}$, which is known to be positive\footnote{We note that the sign of
    $\Delta m_{21}^{2}$ can be known since oscillations in matter are also
    relevant for the solar squared mass difference, which depart from the simple
    formula of Eq.~\eqref{eq:two-flavour-osc}.} (\textit{i.e.}
  $\Delta m^{2}_{21} > 0$). The solar mixing angle $\theta_{\text{sol}}$ is then
  associated with $\theta_{12}$. The atmospheric mixing parameters are
  identified with $|\Delta m_{31}^{2}|$ or $|\Delta m_{32}^{2}|$ and
  $\theta_{23}$. Of course, three-flavour effects alter the simplistic picture
  presented here and must be included to interpret measurements of $\theta_{13}$
  and $\delta_{\text{CP}}$, see \textit{e.g.} Ref.~\cite{Giganti:2017fhf} and
  references therein for a more detailed discussion. The picture that emerges
  from these experiments is then
  \begin{equation}
    \Delta m^{2}_{\text{sol}} \approx \Delta m^{2}_{21} \ll |\Delta m_{31}^{2}| \approx |\Delta m_{32}^{2}| \approx |\Delta m_{\text{atm}}^{2}| \ ,
  \end{equation}
  with which both the \textit{normal mass ordering} $m_{1} < m_{2} < m_{3}$ and
  the \textit{inverted mass ordering} $m_{3} < m_{1} < m_{2}$ are consistent.

  Neutrino oscillation experiments have continued to probe the squared mass
  differences and mixing parameters with impressively high precision, see
  \textit{e.g.} Refs.~\cite{Capozzi:2017ipn, Esteban:2020cvm}. Reactor neutrino
  experiments like KamLAND~\cite{Eguchi:2002dm} and long-baseline accelerators,
  \textit{e.g.} T2K~\cite{Abe:2015awa} and NO$\nu$A~\cite{Adamson:2017qqn}, are
  sensitive to all of these parameters, although $\Delta m^{2}_{\text{atm}}$ and
  $\theta_{13}$ are best measured at short-baseline reactors like Double
  Chooz~\cite{Ardellier:2006mn}, RENO~\cite{Ahn:2010vy} and Daya
  Bay~\cite{An:2015qga}. Today the octant of $\theta_{12}$ is certainly known,
  while $\theta_{13}$ is constrained to be close to 0.15. The sign of the
  atmospheric squared mass difference is still unknown, and therefore so is the
  mass ordering for the neutrinos. The value of the CP-violating Dirac phase
  $\delta_{\text{CP}}$ is less clear, although there is a preference for a value
  somewhere between $\pi$ and $3\pi/2$. The extent of CP-violation in the
  neutrino sector can be represented in a rephasing-invariant way with the
  leptonic Jarlskog invariant
  \begin{equation}
    J_{CP} = \frac{1}{8} \cos \theta_{13} \sin 2\theta_{12} \sin 2\theta_{13} \sin 2\theta_{23} \sin \delta_{\text{CP}} \ ,
  \end{equation}
  so a value of $\pi$ would imply no CP-violating effects, while
  $\delta_{\text{CP}} = 3\pi/2$ would make these maximal. For this work we take
  the results of the most recent fit to neutrino oscillation data by the NuFit
  collaboration~\cite{Esteban:2020cvm, nufitweb} in the context of the
  three-flavour paradigm. These results are summarised in
  Fig.~\ref{fig:nufit-results} separately for the cases of normal and inverted
  mass ordering. Results including atmospheric neutrino oscillation data from
  Super-Kamiokande and those not are also distinguished. Two-dimensional
  projections of the $\chi^{2}$ fit for the same parameters are shown in
  Fig.~\ref{fig:nufit-2d-results}. These data suggest a leptonic mixing matrix
  that has a very different form to the CKM matrix, which we call $\mathbf{V}$.
  We represent this qualitatively by using boxes whose side lengths are scaled
  to the magnitude of the best-fit values of the parameters in the matrices, the
  textures are
  \begin{equation}
    \label{eq:pmns-ckm-comparison}
    \mathbf{U} \sim
    \begin{pmatrix}
       \rule{0.825146em}{0.825146em} & \rule{0.544911em}{0.544911em} & \rule{0.149018em}{0.149018em} \\
       \rule{0.271967em}{0.271967em} & \rule{0.607533em}{0.607533em} & \rule{0.746283em}{0.746283em} \\
       \rule{0.440837em}{0.440837em} & \rule{0.577907em}{0.577907em} & \rule{0.648734em}{0.648734em} \\
     \end{pmatrix}, \qquad
     \mathbf{V} \sim
     \begin{pmatrix}
       \rule{0.97427em}{0.97427em}  & \rule{0.22534em}{0.22534em} & \rule{0.00351em}{0.00351em} \\
       \rule{0.22520em}{0.22520em}  & \rule{0.97344em}{0.97344em} & \rule{0.0412em}{0.0412em} \\
       \rule{0.00867em}{0.00867em}  & \rule{0.0404em}{0.0404em} & \rule{0.999146em}{0.999146em} \\
     \end{pmatrix} \ .
  \end{equation}
  For the PMNS matrix we take the best fit values for the normal mass ordering
  including the Super-Kamiokande results, \textit{i.e.} the numbers in the
  top-left quadrant of Fig.~\ref{fig:nufit-results}. The same numbers also imply
  an upper bound on the mass of the heaviest neutrino at
  \begin{equation}
    \label{eq:atmospheric-bound}
    m_{\text{heaviest}} \leq \sqrt{|\Delta m_{\text{atm}}^{2}|} \approx \SI{0.05}{\eV} \ .
  \end{equation}

  \begin{figure}
    \centering
    \includegraphics[width=0.9\linewidth]{v50.tbl-parameters}
    \caption{The figure shows a table taken from the latest global fit to
      neutrino mass and mixing parameters by the NuFit
      collaboration~\cite{Esteban:2020cvm, nufitweb} in the three-flavour
      picture. The results presented in the upper (lower) panel are obtained by
      excluding (including) the $\chi^{2}$ data on atmospheric neutrinos
      provided by the Super-Kamiokande collaboration (SK). The numbers in the
      1st (2nd) column are obtained assuming normal (inverted) neutrino mass
      ordering. See Ref.~\cite{nufitweb} for more information.}
    \label{fig:nufit-results}
  \end{figure}
 
  \begin{figure}
    \centering
    \includegraphics[width=0.8\linewidth]{v50.fig-region-glob}
    \caption{The figure shows the two-dimensional allowed regions obtained by
      the latest fit to the neutrino mass and mixing parameters by the NuFit
      collaboration~\cite{Esteban:2020cvm, nufitweb}. Each plot shows the
      two-dimensional projection of the allowed region after marginalising with
      respect to the other parameters. The coloured regions (black contour
      curves) are obtained by excluding (including) the Super-Kamiomande
      $\chi^{2}$ data. The different contours correspond to the two-dimensional
      allowed regions at $1\sigma$, $90\%$, $2\sigma$, $99\%$, $3\sigma$
      confidence. See Ref.~\cite{nufitweb} for more information.}
    \label{fig:nufit-2d-results}
  \end{figure}

\subsection{Other experimental probes}

Although neutrino oscillations provide a wealth of evidence for non-zero masses
for at least two of the neutrino fields, they do not probe the absolute mass
scale. There are however kinematic and cosmological probes which bound the
neutrino masses, and some of these are mentioned here.

\subsubsection{Beta decay}

A study of the kinematics of beta-decay experiments shows that differences in
the energy distribution of the emitted electron are expected for different
values of the neutrino mass. Currently these experiments only provide upper
bounds on the effective neutrino mass~\cite{Shrock:1980vy}
\begin{equation}
  m_{\beta} \equiv \sqrt{\sum_{i} |U_{1}^{\ i}|^{2} m_{i}^{2}} \ .
\end{equation}
The best results come from tritium experiments which probe
${}^{3}\mathrm{H} \to {}^{3}\mathrm{He} + e + \bar{\nu}_{e}$. The KATRIN
experiment recently presented the most stringent upper bound~\cite{Aker:2019uuj,
  Aker:2019qfn}
\begin{equation}
  \label{eq:2}
  m_{\beta} < \SI{1.1}{\eV} \ ,
\end{equation}
at $90\%$ confidence, with a projected limit of $m_{\beta} < \SI{0.2}{\eV}$ with
the full dataset.

Another process with the potential to probe the absolute scale of the neutrino
masses, as well as the Majorana phases, is neutrinoless double beta decay
($0\nu\beta\beta$). The process requires the violation of lepton-number by two
units, and is therefore intimately tied to the neutrinos' possible Majorana
nature. If double-beta decay is seen, the \textit{black box
  theorem}~\cite{Schechter:1981bd, Takasugi:1984xr, Hirsch:2006yk} guarantees
that the neutrinos pick up a radiative Majorana mass, even if the neutrino
masses vanish at tree level. Graphically this is easy to understand: a
$0\nu\beta\beta$ operator can always be turned into a neutrino self-energy graph
with four loops. The amplitude for the process is proportional to
\begin{equation}
  \langle m_{\nu} \rangle_{\beta\beta} = \sum_{i} (U_{1}^{\ i})^{2} m_{i} \ ,
\end{equation}
which features both the neutrino masses and the Majorana phases. (Of course it
may be that the four-loop contribution to the neutrino mass implied by the
double-beta-decay operator represents only a small correction to the neutrino
masses, which could arise at lower order.) Current limits on
$\langle m_{\nu} \rangle_{\beta\beta}$ are around \SI{0.2}{\eV}, see
\textit{e.g.} Ref.~\cite{Dolinski:2019nrj} for a recent review of the
experimental status and prospects. Constraints on
$\langle m_{\nu} \rangle_{\beta\beta}$ are usually presented against the mass of
the lightest neutrino mass eigenstate. The behaviour of
$\langle m_{\nu} \rangle_{\beta\beta}$ is very sensitive to the neutrino mass
ordering, in such a way that the inverted scenario implies a minimum allowed
value of $\langle m_{\nu} \rangle_{\beta\beta}$, which will begin to be probed
by the next generation of experiments. A combined global limit and the projected
sensitivity of the DARWIN experiment~\cite{Agostini:2020adk} are shown in
Fig.~\ref{fig:0vbb-darwin}, which also illustrates the different behaviour of
the inverted and normal neutrino mass orderings in the
$\langle m_{\nu} \rangle_{\beta\beta}$ vs. $m_{\text{lightest}}$ plane.

\begin{figure}[t]
  \centering
  \includegraphics[width=0.6\linewidth]{0vbb_darwin}
  \caption{The figure shows limits on the effective neutrino mass for different
    values of $m_{\text{lightest}}$ for both normal and inverted mass ordering.
    The grey region represents the combined sensitivity from a number of leading
    experiments~\cite{Agostini:2019hzm}. The yellow regions are projections for
    the DARWIN experiment under different background hypotheses. The figure is
    taken from Ref.~\cite{Agostini:2020adk}, and we point the reader there for
    more information.}
  \label{fig:0vbb-darwin}
\end{figure}

\subsubsection{Cosmological limits}

The most stringent limits on the sum of the neutrino masses come from cosmology,
although they are model-dependent. In the minimal $\Lambda$CDM model adjusted
for massive neutrinos, the limit implied by the most recent Planck data
release~\cite{Aghanim:2018eyx} is
\begin{equation}
  \sum_{i} m_{i} < \SI{0.12}{\eV} \ ,
\end{equation}
at $95\%$ confidence. This is impressively small, and puts pressure on the
inverted-ordering scenario, for which $\sum_{i} m_{i} \gtrsim \SI{0.1}{\eV}$.
Excitingly, future cosmological probes will likely make a measurement of the sum
of the neutrino masses.

\subsection{Models of neutrino masses}
\label{sec:models-of-mv}

In the SM the neutrino fields appear only within the lepton doublet $L$, and one
cannot write down---in analogy to the up-type Yukawa---a renormalisable operator
that leads to neutrino masses because of the absence of the right-handed fields.
A simple model of neutrino mass then involves introducing right-handed neutrino
fields $\bar{\nu} \sim (\mathbf{1}, \mathbf{1}, 0)_{(\mathbf{1}, \mathbf{2})}$,
extending the Yukawa sector of the SM accordingly to
\begin{equation}
  -\mathscr{L}_{Y} \supset y_{e}\bar{e}LH^{\dagger} + y_{d}\bar{d}QH^{\dagger} + y_{u} \bar{u}QH + y_{\nu}\bar{\nu}LH \ ,
\end{equation}
in a simplified one-flavour picture. This implies Dirac neutrinos with a mass
$m_{\nu} \approx y_{\nu} v$, and makes mass generation symmetric between the
quarks and leptons.

There are a number of problems with this simplistic scenario. First, it is
perhaps uncomfortable to suppose that the neutrinos are many orders of magnitude
lighter than the charged fermions only because $y_{\nu}$ is a very small number.
Indeed, the posited large hierarchy between $y_{\nu}$ and $y_{u}$ seems to spoil
any aesthetic arguments for quark--lepton symmetry in favour of this hypothesis.
Although the Yukawa couplings for the other SM fermions span six orders of
magnitude, couplings within any one generation are all of similar order. We
illustrate this in Fig.~\ref{fig:fermion-masses}, where the fermion masses are
plotted by generation. Whereas one could consider that some underlying theory of
flavour may explain, for example, the disparity in scale between the masses of
the electron and the top quark, the large mass difference between the electron
and the lightest neutrino, although technically natural, seems to cry out for
its own explanation.

A second point of criticism with the simple scenario presented above is that it
ignores the Majorana mass term for $\bar{\nu}$
\begin{equation}
  \label{eq:3}
  \mathscr{L} \supset - \tfrac{1}{2} \mu \bar{\nu} \bar{\nu} + \text{h.c.}
\end{equation}
In order to maintain the Dirac neutrino mass with the relation
$m_{\nu} \approx y_{\nu}v$, the Majorana mass must be forbidden or else chosen
to be negligibly small, assumptions adding additional layers of contrivance to
the theory. A sensible choice for the scale $\mu$ is some high scale $\Lambda$
associated with the breaking of $\mathrm{U}(1)_{B-L}$. Taking
$\mu \gg y_{\nu} v$, the neutrino-mass matrix takes on the form
\begin{equation}
  \mathbf{m}_{\nu} = \begin{pmatrix}
    0 & y_{\nu} v \\
    y_{\nu} v & \mu \\
  \end{pmatrix} \ ,
\end{equation}
with eigenvalues
\begin{equation}
  m_{l} \approx \frac{y_{\nu}^{2}v^{2}}{\mu},\quad m_{h} \approx \mu \ .
\end{equation}
The assumption $\mu \gg y_{\nu} v$ implies that $m_{l} \ll y_{\nu} v$, and the
neutrino is successfully arranged to be much lighter than $y_{\nu} v$, which can
now be taken to be on the order of the charged fermion masses. The theory also
leaves us with a neutrino whose mass must be significantly larger than the
electroweak scale: $m_{h} \gtrsim \SI{e14}{\GeV}$, assuming
$m_{l} = \SI{0.05}{\eV}$ and $y_{\nu} = 1$. After transforming into the
mass-diagonal basis, the physical fields $\nu_{l}$ and $\nu_{h}$ correspond to
Majorana particles. Thus, in the most motivated region of parameter space, the
phenomenology of the light neutrinos in even the $\text{SM} + \bar{\nu}$
scenario is Majorana. This toy scenario illustrates the mechanism commonly
called the \textit{seesaw mechanism}: making the neutrinos very light at the
expense of making other fields very heavy. This is discussed more broadly below.

\begin{figure}[t]
  \centering
  \includegraphics[width=0.7\linewidth]{fermion_mass_plot}
  \caption{The masses of the SM fermions grouped by generation. While the SM
    Yukawa couplings span a wide range of values, within any specific generation
    they are all of similar order. The tiny masses of the neutrinos seem to
    suggest an alternate mass-generation mechanism.}
  \label{fig:fermion-masses}
\end{figure}

\subsubsection{Tree-level models of neutrino mass}

The toy seesaw scenario discussed above can be understood more generally by
studying the effective theory valid below the scale $\mu$, which does not
contain the field $\bar{\nu}$. The leading-order lepton-number-violating (LNV)
effects appear at dimension five in the operator
\begin{equation}
  \mathscr{L} \supset \frac{\kappa}{\Lambda} (L^{i}L^{j})H^{k}H^{l} \cdot \epsilon_{ik}\epsilon_{jl} \ ,
\end{equation}
with $\kappa$ a dimensionless coefficient. This operator is commonly called the
\textit{Weinberg operator}. In the broken phase it gives rise to a Majorana mass
for the neutrinos consistent with the seesaw formula:
\begin{equation}
  \label{eq:seesaw-formula}
  \mathscr{L} \supset \frac{v^{2} \kappa}{\Lambda} \nu \nu \ .
\end{equation}

The $\text{SM}+\bar{\nu}$ scenario is not the only simplified model that
generates the Weinberg operator at tree-level. A simple diagram-topology
analysis suggests there are another two seesaw mechanisms in the UV: a model
with a scalar isotriplet field $\Xi_{1} \sim (\mathbf{1}, \mathbf{3}, 1)_{S}$,
called the type-II mechanism, and another with an isotriplet Majorana fermion
$\Sigma \sim (\mathbf{1}, \mathbf{3}, 0)_{F}$, called type-III. Along with the
type-I heavy $\bar{\nu}$ model, these are collectively referred to as the
canonical seesaw mechanisms~\cite{MINKOWSKI1977421, Yanagida:1979as,
  GellMann:1980vs, PhysRevLett.44.912, Glashow:1979nm, Magg:1980ut,
  PhysRevD.22.2227, LAZARIDES1981287, Wetterich:1981bx, PhysRevD.23.165,
  Foot:1988aq} and have been studied at length in the literature. They are
simple models in that they introduce only very few exotic degrees of freedom and
free parameters. However, the high seesaw scale makes these models practically
untestable at current and future collider experiments.

Models of Majorana neutrino mass can be made more testable if, instead of the
suppression of the neutrino masses coming from a large $\Lambda$ in
Eq.~\eqref{eq:seesaw-formula}, the coefficient $\kappa$ were somehow arranged to
be small. Although there are a number of mechanisms to achieve this, we
concentrate below on radiative models of Majorana neutrino mass, in which
$\kappa$ is made small through loop and coupling suppression.

\subsubsection{Radiative models and their classification}

It may be the case that the field content of whatever high-energy theory
describes the neutrino masses is such that no neutrino self-energy diagram can
be drawn at the tree level. Indeed, this will be the case if there is
lepton-number violation by two units ($\Delta L = 2$) from interactions other
than those present in the canonical seesaw models. Such models are called
radiative, since the neutrino masses arise through loop graphs. The historically
important Zee~\cite{Zee:1980ai} and Zee--Babu~\cite{Zee:1985id, Babu:1988ki}
models have come to be archetypal radiative scenarios in which the neutrinos
gain masses through $\Delta L = 2$ the interactions of exotic scalars at one and
two loops respectively. In the Zee model, an additional Higgs doublet
$\varphi \sim (\mathbf{1}, \mathbf{2}, \tfrac{1}{2})$ and a charged scalar
$\mathcal{S}_{1} \sim (\mathbf{1}, \mathbf{1}, 1)$ are introduced, while the
Zee--Babu model contains $\mathcal{S}_{1}$ and the doubly-charged scalar
$\mathcal{S}_{2} \sim (\mathbf{1}, \mathbf{1}, 2)$. The neutrino self-energy
diagrams for these models are shown in Fig.~\ref{fig:zee-and-zee-babu-diags} as
examples. The corresponding neutrino-mass matrices are
\begin{align}
  [\mathbf{m}^{\text{Zee}}_\nu]_{rs} & = \frac{\mu v^2}{16\pi^2} \sum_t x_{[rt]} I_{t} [\mathbf{m}_e]_t y_{ts} + (r \leftrightarrow s) \ , \label{eq:zee-mv} \\
  [\mathbf{m}^{\text{Zee--Babu}}_\nu]_{rs} & = \frac{\mu^{\prime} v^2}{(16 \pi^2)^2} \sum_{t,u} x_{[rt]} [\mathbf{m}_e]_t z_{\{tu\}} [\mathbf{m}_e]_u x_{[us]} I^\prime_{tu} + (r \leftrightarrow s) \ , \label{eq:zee-babu-mv}
\end{align}
written in terms of the couplings $x,y,z \in \mathbb{C}$ as shown in
Fig.~\ref{fig:zee-and-zee-babu-diags}, and the associated loop functions $I$ and
$I^{\prime}$. One can see that in these models, the coefficient of the Weinberg
operator $\kappa$ is naturally suppressed with respect to the seesaw formula by
SM-fermion masses (charged-lepton masses in this case), exotic couplings, and
loop factors.

Such models are economic, since they do not require the imposition of \textit{ad
  hoc} symmetries, and in many cases make a connection to other unsolved
problems of the SM such as the nature of dark matter or the matter--antimatter
asymmetry of the Universe. They are also elegant, since the smallness of the
neutrino masses emerges as a natural consequence, rather than through the
imposed requirement of exceedingly small coupling constants. Radiative models
are easier to probe experimentally since the additional loop suppression and
products of couplings bring down the allowed scale of the new physics, in some
cases to LHC-accessible energy ranges~\cite{deGouvea:2007qla}. The Zee--Babu
model, for example, is non-trivially constrained by same-sign dilepton searches
performed by ATLAS~\cite{ATLAS:2012hi, ATLAS:2014kca, Aaboud:2017qph} and
CMS~\cite{Chatrchyan:2012ya, CMS:2016cpz, CMS:2017pet}. Additionally, the
predicted dependence of the neutrino-mass matrix on SM parameters, as is clear
from Eqs.~\eqref{eq:zee-mv} and \eqref{eq:zee-babu-mv}, can make these models
very predictive as far as neutrino phenomenology is concerned. For example, the
minimal version of the Zee model has the charged-lepton masses arising from
couplings to $\varphi$, allowing for a simplification of the expression
Eq.~\eqref{eq:zee-mv} such that the leptonic flavour index $t$ is identified
with $s$. The antisymmetry in the leptonic flavour indices implied by $x_{[rs]}$
must be accounted for by an antisymmetry in the loop integral such that
$I \sim I_{r} - I_{s}$. The flavour structure then
becomes~\cite{Wolfenstein:1980sy}
\begin{equation}
  \label{eq:zee-wolfenstein-mv}
  [\mathbf{m}_{\nu}^{\text{Min. Zee}}]_{rs} \propto x_{[rs]} ([\mathbf{m}_{e}]_{r}^{2} - [\mathbf{m}_{e}]_{s}^{2}) \ ,
\end{equation}
which makes clear predictions relevant to neutrino oscillations that in this
case rule out the model. For recent reviews of radiative models see
Refs.~\cite{Boucenna:2014zba, Cai:2017jrq}.

\begin{figure}[t]
  \centering
  \includegraphics[width=0.45\textwidth]{zee_diag}
  \includegraphics[width=0.45\textwidth]{zee_babu_diag}
  \caption{The neutrino self-energy diagrams relevant to the Zee (left) and
    Zee--Babu (right) models.}
  \label{fig:zee-and-zee-babu-diags}
\end{figure}

The Zee and Zee--Babu scenarios are only two of a very large number of radiative
models, none of which are \textit{a priori} more likely to be true than any
other. In the context of such a large theory-space, it is useful to have an
organising principle to aid in the study and classification of these models, and
beginning with $\Delta L = 2$ effective operators has been shown to be an
effective strategy.

One approach to this model taxonomy involves studying loop-level completions of
the Weinberg operator, and its dimension-($5+2n$) generalisations
\begin{equation*}
 \mathcal{O}_1^{\prime \cdots \prime} = (LL)HH(H^\dagger H)^n  \ .
\end{equation*}
Here, models can be systematically written down by studying the various
topologies able to be accommodated by the operator with increasing number of
loops. This is done in such a way that models implying lower-order contributions
to the neutrino mass can be discarded~\cite{Farzan:2012ev}. Such an approach has
been applied to the Weinberg operator up to three loops~\cite{Bonnet:2012kz,
  Sierra:2014rxa, Cepedello:2018rfh} and to its dimension-seven generalisation
at one loop~\cite{Cepedello:2017eqf}. An alternative and complementary method
begins by considering all of the gauge-invariant $\Delta L = 2$ operators in the
SM effective field theory (SMEFT), first listed in this context by Babu and
Leung (BL)~\cite{Babu:2001ex} and extended by de Gouv\^{e}a and Jenkins
(dGJ)~\cite{deGouvea:2007qla}. Supposing that the tree-level coefficient of one
of these is non-zero at the high scale, neutrino masses will be generated from
loop graphs contributing to the mixing of this operator and the Weinberg-like
operators $\mathcal{O}_1^{\prime \cdots \prime}$. The process of expanding the
operator into a series of UV-complete, renormalisable models that generate the
parent operator at tree-level is called \emph{opening up} or \emph{exploding}
the operator. The remaining external fields must be looped-off, with additional
loops of SM gauge bosons or Higgs fields added as necessary in order to obtain a
neutrino self-energy diagram. A model-building formula along these lines has
been formulated in Ref.~\cite{PhysRevD.87.073007}, and it has been used to write
down all of the minimal, tree-level UV-completions of $\Delta L = 2$ operators
at dimension seven~\cite{Cai:2014kra} corresponding to tree-level and radiative
neutrino-mass models. The tree-level completions of the Weinberg-like operators
have been written down up to dimension eleven~\cite{Cai:2014kra, Bonnet:2009ej,
  Anamiati:2018cuq}. The operators $\Delta L = 2$ operators in the SM EFT are
discussed in more detail in Sec.~\ref{sec:EFT-of-SM}.

We note that an economic classification scheme, separate from an EFT framework,
was presented in Ref.~\cite{Klein:2019iws} based on the number of exotic degrees
of freedom by which the SM is extended. There, the method is applied to the case
of radiative models with two exotics\footnote{Including models with one scalar
  and one Dirac fermion.}, and has also been used to study minimal neutrino-mass
models compatible with $\mathrm{SU}(5)$ unification~\cite{Klein:2019jgb}.

\section{Effective field theories of the SM}

In the following we introduce EFTs in general, and then specialise to those
built out of SM fields: the Standard Model EFT (SMEFT) and the Weak Effective
Theory (WET). We place particular emphasis on the operators appearing at
dimension-six, and then at odd mass dimension up to dimension eleven. The
odd-dimension operators in the SMEFT are those that organise the building of
neutrino-mass models in the tradition of Refs.~\cite{Babu:2001ex,
  deGouvea:2007qla, PhysRevD.87.073007, Cai:2014kra}. We begin with prefatory
comments on the process of tree-level matching, then move on to discuss the
SMEFT and peripheral concepts.

\subsection{Tree-level matching}

Suppose one has a theory with light particle states described by fields
$\pi_{i}$ and heavy states described by $\Pi_{i}$ with a Lagrangian of the form
\begin{equation}
  \begin{aligned}
  \mathscr{L}_{\text{UV}}[\pi, \Pi] &= \mathscr{L}_{\text{kin}}[\pi, \Pi] + \mathscr{L}_{\text{int}}[\pi, \Pi], \text{ with } \\
  \mathscr{L}_{\text{int}}[\pi, \Pi] &= \mathscr{L}^{l}[\pi] + \mathscr{L}^{h}[\Pi] +  \mathscr{L}^{lh}[\pi, \Pi].
  \end{aligned}
\end{equation}
Below the threshold for $\Pi_{i}$ production, an effective description of the
theory can be used that involves interactions only between the light fields.
This effective theory is described by a Lagrangian
$\mathscr{L}_{\text{eff}}[\pi]$ involving interactions between the
$\pi_{i}$ that correspond to diagrams in the full theory containing only heavy
internal propagators and light external states. At the classical level,
$\mathscr{L}_{\text{eff}}$ can be written down by integrating out the
$\Pi_{i}$. Perturbatively this corresponds to expanding the heavy propagators
$\Delta$ in powers of momenta on the heavy mass scale\footnote{We note that some
  UV scenarios may have more than one characteristic scale. In this case
  $\Lambda$ can be understood as an effective scale which may not necessarily
  correspond to the mass of a specific particle.} $\Lambda$, such
that
\begin{equation}
  \label{eq:prop}
  \Delta = \begin{dcases}
    -\frac{1}{\Lambda^{2}}\left(1 + \frac{p^2}{\Lambda^{2}} + \cdots \right) & \text{for
    }\quad \includegraphics[width=0.17\linewidth]{scalar_prop.pdf} \\
    -\frac{\delta_{\alpha}^{\ \beta}}{\Lambda} \left( 1 + \frac{p^2}{\Lambda^{2}} + \cdots \right) & \text{for
    }\quad \includegraphics[width=0.17\linewidth]{fermion_mprop.pdf} \\
    -\frac{i p \cdot \bar{\sigma}^{\dot{\alpha}\beta}}{\Lambda^{2}} \left( 1 + \frac{p^2}{\Lambda^{2}} + \cdots \right) & \text{for
    }\quad \includegraphics[width=0.17\linewidth]{fermion_pprop.pdf}
  \end{dcases}
  \ .
\end{equation}
In this notation, the arrow-preserving propagator corresponds to the part of the
regular four-component fermion propagator proportional to momentum, while the
arrow-violating one is the part proportional to the fermion mass. Expressions
for the fermion propagators with reversed arrows follow from
$\bar{\sigma}^{\mu} \to \sigma^{\mu}$ and interchanging dotted and undotted
indices (see Ref.~\cite{Dreiner:2008tw} Sec.~4.2 for the Lorentz structure).

Equivalently, the integration can be performed using the classical EOM of the
$\Pi_{i}$. For some heavy field $\Pi$, the linearised solution to its classical
EOM can be used to remove it from the Lagrangian completely. This procedure is
mildly different for scalars and fermions, and we briefly outline these
separately below. In both cases, we begin with a Lagrangian
$\mathscr{L}_{\text{UV}}$ for which we imagine that kinetic and mass mixing
terms between heavy and light fields have been removed.

There are tree-level contributions to $\mathscr{L}^{\text{eff}}$ as long as
there are interaction terms linear in $\Pi$. For scalar $\Pi$, the UV Lagrangian
contains the terms
\begin{equation}
  \label{eq:scalar-lag}
 \mathscr{L}_{\text{UV}}[\Pi, \pi] \supset \Pi^{\dagger} (- D^{2} - m_{\Pi}^{2}) \Pi + \left(\Pi \frac{\partial \mathscr{L}^{lh}}{\partial \Pi} + \text{h.c.} \right) \ ,
\end{equation}
where $\partial \mathscr{L}^{lh} / \partial \Pi$ is a function only of
light fields, and we are neglecting interactions of the form
$\Pi^{\dagger} \Pi f(\pi)$ for the sake of conciseness. The EOM are
\begin{equation}
  \label{eq:scalar-eom}
 (- D^{2} - m_{\Pi}^{2}) \Pi = - \frac{\partial \mathscr{L}^{lh}}{\partial \Pi^{\dagger}} + \mathcal{O}(\Pi^{2}) \ ,
\end{equation}
which can be solved for $\Pi^{\text{cl}}$, the classical field configuration, by
inverting the differential operator on the LHS of Eq.~\eqref{eq:scalar-eom} and
expanding in $D^{2} / m_{\Pi}^{2}$:
\begin{equation}
  \label{eq:scalar-repl}
  \Pi^{\text{cl}} = - \frac{1}{m_{\Pi}^{2}} \left( 1 - \frac{D^{2}}{m_{\Pi}^{2}}  + \cdots \right) \frac{\partial \mathscr{L}^{lh}}{\partial \Pi^{\dagger}} \ .
\end{equation}
This solution can be substituted back into Eq.~\eqref{eq:scalar-lag} to give
interactions between light fields in the tree-level effective Lagrangian:
\begin{equation}
  \label{eq:classical-efflag-scalar}
  \mathscr{L}_{\text{eff}}[\pi] \supset - \frac{\partial \mathscr{L}^{lh}}{\partial \Pi} \frac{1}{m_{\Pi}^{2}} \left( 1 - \frac{D^{2}}{m_{\Pi}^{2}} + \cdots \right) \frac{\partial \mathscr{L}^{lh}}{\partial \Pi^{\dagger}} \ .
\end{equation}
Many concrete examples of this procedure can be found in the literature, see
\textit{e.g.} Ref.~\cite{Henning:2014wua}. The expansion in $D^{2}/m_{\Pi}^{2}$
corresponds to the expansion in $p^{2} / \Lambda^{2}$ in the first case of
Eq.~\eqref{eq:prop}, showing the expansion of the scalar propagator.

Next we sketch out the procedure for a Dirac fermion
$\Pi + \bar{\Pi}^{\dagger}$, where $\Pi$ and $\bar{\Pi}$ are separate
two-component spin-$\tfrac{1}{2}$ fields transforming oppositely under
$G_{\text{SM}}$. In this case, the UV theory is described by a Lagrangian like
\begin{equation}
  \label{eq:fermion-lag}
  \mathscr{L}_{\text{UV}}[\Pi, \pi] \supset i \Pi^{\dagger} \bar{\sigma}^{\mu} D_{\mu} \Pi + i \bar{\Pi}^{\dagger} \bar{\sigma}^{\mu} D_{\mu} \bar{\Pi} + \left( \Pi \frac{\partial \mathscr{L}^{lh}}{\partial \Pi} + \bar{\Pi} \frac{\partial \mathscr{L}^{lh}}{\partial \bar{\Pi}} - m_{\Pi} \bar{\Pi} \Pi  + \text{h.c.} \right)
\end{equation}
Varying the action with respect to the heavy fields gives two coupled EOM:
\begin{align}
 i \bar{\sigma}^{\mu} D_{\mu} \Pi - m \bar{\Pi}^{\dagger} + \frac{\partial \mathscr{L}^{lh}}{\partial \Pi^{\dagger}} &= 0 \ , \label{eq:pi-fermion-eom} \\
 i \bar{\sigma}^{\mu} D_{\mu} \bar{\Pi} - m \Pi^{\dagger} + \frac{\partial \mathscr{L}^{lh}}{\partial \bar{\Pi}^{\dagger}} &= 0 \ . \label{eq:pibar-fermion-eom}
\end{align}
Substituting Eq.~\eqref{eq:pi-fermion-eom} into Eq.~\eqref{eq:pibar-fermion-eom}
gives a second-order partial differential equation in $\Pi$, analogous to
Eq.~\eqref{eq:scalar-eom}. Inverting the differential operator in a similar way
gives
\begin{equation}
  \label{eq:fermion-repl}
  \Pi^{\text{cl}}_{\beta} = \frac{1}{m_{\Pi}^{2}} \left( \epsilon_{\alpha \beta} + \frac{ \tfrac{1}{2} X_{\alpha \beta} - D^{2} \epsilon_{\alpha \beta}}{m_{\Pi}^{2}} + \cdots \right) \left( i D^{\alpha \dot{\alpha}} \frac{\partial \mathscr{L}^{lh}}{\partial \Pi^{\dagger}_{\dot{\beta}}} \epsilon_{\dot{\alpha} \dot{\beta}} + m_{\Pi} \frac{\partial \mathscr{L}^{lh}}{\partial \bar{\Pi}_{\alpha}} \right) \ ,
\end{equation}
where the field-strength tensor comes about from a structure like
\begin{align}
  [\sigma^{\mu} \bar{\sigma}^{\nu}]_{\alpha}^{\ \beta} D_{\mu} D_{\nu} &= \eta^{\mu\nu} D_{\mu} D_{\nu} \delta_{\alpha}^{\ \beta} - 2i [\sigma^{\mu\nu}]_{\alpha}^{\ \beta} D_{\mu} D_{\nu} \\
  &= D^{2} \delta_{\alpha}^{\ \beta} - \tfrac{1}{2} X_{\alpha}^{\ \beta} \ .
\end{align}
Here, and later in this section, the replacement $\bar{\Pi} \to \Pi$ should be
understood for Majorana $\Pi$. Each contribution corresponds to a particular
kind of propagator in the perturbative picture. The first term in the last
parenthesis of Eq.~\eqref{eq:fermion-repl} results from the fermion propagator
proportional to momentum: the arrow-preserving fermion propagator shown as the
last case of Eq.~\eqref{eq:prop}. The second term in the same parentheses stems
from the fermion propagator proportional to the mass, corresponding to the
arrow-violating propagator shown in the middle case of Eq.~\eqref{eq:prop}.
Replacing $\Pi$ in Eq.~\eqref{eq:fermion-lag} gives the tree-level effective
Lagrangian with the heavy fermion integrated out:
\begin{equation}
  \begin{aligned}
  \label{eq:classical-efflag-fermion}
  \mathscr{L}_{\text{eff}}[\pi] &\supset \frac{\partial \mathscr{L}^{lh}}{\partial \Pi_{\beta}} \frac{1}{m_{\Pi}^{2}} \left(\epsilon_{\alpha \beta} + \frac{ \tfrac{1}{2} X_{\alpha \beta} - D^{2} \epsilon_{\alpha \beta}}{m_{\Pi}^{2}} +  \cdots \right) i D^{\alpha \dot{\alpha}} \frac{\partial \mathscr{L}^{lh}}{\partial \Pi^{\dagger}_{\dot{\beta}}} \epsilon_{\dot{\alpha} \dot{\beta}}\\
  &\quad + \frac{\partial \mathscr{L}^{lh}}{\partial \Pi_{\beta}} \frac{1}{m_{\Pi}^{2}} \left(\epsilon_{\alpha \beta} + \frac{ \tfrac{1}{2} X_{\alpha \beta} - D^{2} \epsilon_{\alpha \beta}}{m_{\Pi}^{2}} + \cdots \right) \frac{\partial \mathscr{L}^{lh}}{\partial \bar{\Pi}_{\alpha}} \ .
  \end{aligned}
\end{equation}

As shown in Eqs.~\eqref{eq:classical-efflag-scalar} and
\eqref{eq:classical-efflag-fermion}, expanding in powers of derivatives on heavy
masses leads to a tower of local operators of increasing mass dimension $d_{i}$
organised as a power series in the inverse heavy scale:
\begin{equation}
  \mathscr{L}_{\text{eff}}[\pi] = \mathscr{L}^{l}[\pi] + \sum_{i} \frac{C_{i}}{\Lambda^{d_{i}-4}} \mathcal{O}_{i}[\pi].
\end{equation}
The $C_{i}$ are dimensionless coefficients which are in general calculable if
one knows the high-energy theory.

\subsection{Effective field theories of the SM}
\label{sec:EFT-of-SM}

Below we discuss EFTs constructed from SM fields and invariant under SM
symmetries. The main theory of study is the SMEFT: the gauge- and
Lorentz-invariant EFT built from the fields listed in Table~\ref{tab:sm-fields}.
We also mention the WET, also known as the LEFT (Low-energy Effective Field
Theory), for which invariance under
$\mathrm{SU}(2)_{+} \otimes \mathrm{SU}(2)_{-} \otimes \mathrm{SU}(3)_{c} \otimes \mathrm{U}(1)_{\text{EM}}$
is required. In addition to the difficulties associated with constructing the
invariants in an EFT, there are additional intricacies associated with
redundancies among operators. These are also discussed below in the context of
the SMEFT, although the principles discussed apply more generally. We tie off
this section with a discussion of the $\Delta L = 2$ operators in the SMEFT.

\subsubsection{The SMEFT at dimension six}

Given the broad experimental success of the SM, it is perhaps a sensible
assumption that there should be a sizeable mass gap between the electroweak
scale and the mass-scale characterising any new physics. In this context, the
SMEFT can be a powerful tool for constraining what this new physics might look
like in a model-independent way. Indeed, the SMEFT operators at dimension six
are already coming to play an increasingly important role in particle
phenomenology, and they have become the \textit{de facto} framework for
interpreting low-energy constraints on theoretical models, and experimental
deviations from SM predictions. For an extensive review, we point the reader to
Ref.~\cite{Brivio:2017vri}.

It has not been easy to write down a complete basis of operators in the
SMEFT~\cite{}, although this is now a mostly solved problem~\cite{}. The
lowest-dimensional operator appearing in the EFT is also the only dimension-five
operator:
\begin{equation}
  \mathscr{L}^{(5)} = [C_{5}]_{\{rs\}} (L^{i}_{r} L^{j}_{s}) H^{k} H^{l} \epsilon_{ik} \epsilon_{jl} + \text{h.c.} \ ,
\end{equation}
already discussed briefly in Sec.~\ref{sec:models-of-mv}. The matrix of operator
coefficients is necessarily symmetric by Fermi--Dirac statistics. The operator
violates lepton-number by two units, and usually gives the dominant contribution
to the neutrino mass in Majorana models. Ref.~\cite{Kobach:2016ami} shows that
operators in the SMEFT of dimension $d$ satisfy
\begin{equation}
  \frac{1}{2} (\Delta B - \Delta L) = d \mod 2 \ ,
\end{equation}
and thus odd mass-dimension operators must violate $B - L$ by two units, while
operators of even mass-dimension cannot violate $B - L$. So, aside from
lepton-number-violating effects, the leading-order deviations from the SM appear
at dimension six, where there are many more operators.

The dimension-six operators come in eight general classes: $X^{3}$, $H^{6}$,
$H^{4}D^{2}$, $X^{2}H^{2}$, $\psi^{2}H^{3}$, $\psi^{2}XH$, $\psi^{2}H^{2}D$ and
$\psi^{4}$. (Here, $X$ represents a general field-strength tensor, $D$ is a
covariant derivative and $\psi$ is a fermion field.) The most common basis found
in the literature is the Warsaw basis, which tends to prefer not mixing colour
indices in fermion currents.

We do not give a complete listing of these operators here, but instead define
them as necessary in the chapters they appear. For the SMEFT we attempt to
adhere to the conventions of Ref.~\cite{} relevant for the Warsaw basis.

The operators in the notation that we use them throughout the rest of this work
are listed in Table~\ref{tab:smeft-d6-psi4}. Each operator is given in
four-component notation along with the number of operators for $n_{f} = 3$
SM-fermion flavours and the operator class.


\subsubsection{Operator redundancies and the Hilbert series}
\lipsum[2]

\subsubsection{The WET}
\lipsum[2]

\subsubsection{Lepton-number-violating SMEFT}
\lipsum[2]

\section{The flavour anomalies and their explanation}
\lipsum[2]

\subsection{Neutral-current anomalies}
\lipsum[2]

\subsection{Charged-current anomalies}
\lipsum[2]

\subsection{Anomalous magnetic moment of the muon}
\lipsum[2]

\subsection{Scalar and Vector leptoquarks}

\myglossaryentry{lipsum}{lipsum}{Lorem Ipsum, a special type of fudge}{}
\myglossaryentry{dolor}{dolor}{No idea why}{parent={lipsum}}
\myglossaryentry{ibit}{ibit}{Sounds right, doesn't it?}{parent={lipsum}}
\myacronym{DFT}{density functional theory}
\myglossaryentry{$\pi$}{pi}{Greek letter pi, \ensuremath{\Pi} does this work?}{symbol={$\pi$}}
\myacronym{RDF}{radial distribution function}
\myglossaryentry{radial distribution function}{radialdistributionfunction}{}{symbol={$g(r)$}}
